\documentclass{article}
\usepackage[utf8]{inputenc}
\usepackage[T1]{fontenc}
\usepackage{geometry}
\geometry{
    a4paper,
    left=2.5cm,
    right=2.5cm,
    top=2.5cm,
    bottom=2.5cm
}
\usepackage{graphicx}
\usepackage{float}
\usepackage{hyperref}
\usepackage[english]{babel}
\usepackage{amsmath}
\usepackage{amssymb}
\usepackage{array}
\usepackage{booktabs}
\usepackage{caption}
\usepackage{subcaption}
\usepackage{enumitem}
\usepackage{fancyhdr}
\usepackage{titlesec}
\usepackage{url}
\usepackage{textcomp}
\usepackage{xcolor}
\usepackage{listings}
\usepackage{multirow}
\usepackage{longtable}
\usepackage{tabularx}
\usepackage{times}

% Set up page style
\pagestyle{fancy}
\fancyhf{}
\rfoot{\thepage}
\renewcommand{\headrulewidth}{0pt}

% Title formatting
\titleformat{\chapter}[display]{\normalfont\bfseries}{\chaptertitle\thechapter}{0pt}{\Huge}
\titlespacing*{\chapter}{0pt}{50pt}{40pt}

% Document info
\title{Template Document}
\author{Academic Support Agent}
\date{\today}

\begin{document}
\maketitle

%%%%%%%%%%%%%%%%%%%%%%%%%%%%%%%%%%%%%

 %%%%%%%%%%%%%%%%%%%%%%%%%%%%%%%%%%%%%
\maketitle % Dodaj stron  tytu ow 
 %%%%%%%%%%%%%%%%%%%%%%%
 % Je eli chcesz komu  podzi kowa , to mo esz u y  poni szego kodu

\thispagestyle{empty}
\vspace*{\fill}
\begin{flushright}
    \em
    \begin{minipage}{0.75\textwidth}
        Tutaj mo esz umie ci  tre  podzi kowa .
        Tutaj mo esz umie ci  tre  podzi kowa .
        Tutaj mo esz umie ci  tre  podzi kowa .
        Tutaj mo esz umie ci  tre  podzi kowa .
        Tutaj mo esz umie ci  tre  podzi kowa .
    \end{minipage}
\end{flushright}
 %%%%%%%%%%%%%%%%%%%%%%%
\begin{abstractPL}
    Streszczenie po polsku \ldots
\end{abstractPL}
\begin{abstractEN}
    Abstract in English \ldots
\end{abstractEN}
 %%%%%%%%%%%%%%%%%%%%%%%%%%%%%%%%%%%%%
 %%%%%%%%%%%%%%%%%%%%%%%%%%%%%%%%%%%%%

\begin{shaded}
    Zawarto  spisu tre ci \pauza tytu y rozdzia w oraz ich liczba zale  od tematyki pracy \pauza nale y ustali  z opiekunem pracy.
\end{shaded}


\listofalgorithms  % Wygeneruj list  algorytm w
 % Wygeneruj list  kod w  r d owych
\lstlistoflistings % Je eli do tworzenia listing w u ywasz pakietu 'listings'
%  \listoflistings % Je eli do tworzenia listing w u ywasz pakietu 'minted'
\printnomenclature % Wy wietl list  symboli
 %%%%%%%%%%%%%%%%%%%%%%%%%%%%%%%%%%%%%

 %%%%%%%%%%%%%%%%%%%%%%%%%%%%%%%%%%%%%
 % Rozdzia y cz ci g wnej
% Content from tex/introduction would be included here

% Content from tex/sota would be included here

% Content from tex/research would be included here

% Content from tex/conclusions would be included here

 %%%%%%%%%%%%%%%%%%%%%%%%%%%%%%%%%%%%%
 %%%%%%%%%%%%%%%%%%%%%%%%%%%%%%%%%%%%%

 %%%%%%%%%%%%%%%%%%%%%%%%%%%%%%%%%%%%%
 % Rozdzia y dodatku
% Content from tex/example would be included here

 %%%%%%%%%%%%%%%%%%%%%%%%%%%%%%%%%%%%%
 %%%%%%%%%%%%%%%%%%%%%%%%%%%%%%%%%%%%%

 %%%%%%%%%%%%%%%%%%%%%%%%%%%%%%%%%%%%%
 % Rozdzia y cz ci ko cowej
% Content from tex/notes would be included here

 %%%%%%%%%%%%%%%%%%%%%%%%%%%%%%%%%%%%%
 %%%%%%%%%%%%%%%%%%%%%%%%%%%%%%%%%%%%%
 % UWAGA: Je eli b dziesz tworzy  spis literatury r cznie, tj. za pomoc  otoczenia 
 %  'thebibliography', to pami taj --- powinien zawiera  tylko te publikacje, 
 %  do kt rych odwo ujesz si  w pracy. 

\printbibliography  % Wy wietl spis literatury
 %%%%%%%%%%%%%%%%%%%%%%%%%%%%%%%%%%%%%
 %%%%%%%%%%%%%%%%%%%%%%%%%%%%%%%%%%%%%
\end{document}